\documentclass[12pt,oneside]{book}

% ============================================================
% METADATOS
% ============================================================
\newcommand{\universidad}{Universidad de Buenos Aires (UBA)}
\newcommand{\materia}{Derecho Procesal Civil}
\newcommand{\titulo}{Contestación de Demanda y Excepciones Previas}
\newcommand{\autor}{Isaac Edgar Camacho Ocampo}
\newcommand{\anio}{2024-2025}

% ============================================================
% PAQUETES
% ============================================================
\usepackage[utf8]{inputenc}
\usepackage[T1]{fontenc}
\usepackage[provide=*,spanish]{babel}

\usepackage[
    top=2.5cm,
    bottom=2.5cm,
    left=3cm,
    right=2.5cm,
    headheight=15pt
]{geometry}

\usepackage{amsmath}
\usepackage{amssymb}
\usepackage{mathtools}

\usepackage{graphicx}
\usepackage{float}
\usepackage{caption}
\usepackage{subcaption}

\usepackage{booktabs}
\usepackage{array}
\usepackage{longtable}
\usepackage{tabularx}

\usepackage{enumitem}

\usepackage{xcolor}
\usepackage[most]{tcolorbox}

\usepackage{fancyhdr}
\usepackage{ragged2e}

\graphicspath{
    {./img/}
    {./graficos/}
    {./figuras/}
}

\usepackage[
    colorlinks=true,
    linkcolor=blue!50!black,
    citecolor=green!40!black,
    urlcolor=blue!60!black,
    pdftitle={\titulo},
    pdfauthor={\autor},
    pdfsubject={Apuntes universitarios},
    bookmarks=true
]{hyperref}

% ============================================================
% CONFIGURACIÓN GENERAL
% ============================================================
\setcounter{tocdepth}{2}

\setlength{\parindent}{0pt}
\setlength{\parskip}{0.5em}

\setlist[itemize]{
    topsep=0.3em,
    itemsep=0.2em
}

\setlist[enumerate]{
    topsep=0.3em,
    itemsep=0.2em
}

\clubpenalty=10000
\widowpenalty=10000

% ============================================================
% ENCABEZADOS Y PIES
% ============================================================
\pagestyle{fancy}
\fancyhf{}

\fancyhead[L]{\nouppercase{\leftmark}}
\fancyhead[R]{\materia}

\fancyfoot[C]{\thepage}

\renewcommand{\headrulewidth}{0.4pt}
\renewcommand{\footrulewidth}{0pt}

\fancypagestyle{plain}{
    \fancyhf{}
    \fancyfoot[C]{\thepage}
    \renewcommand{\headrulewidth}{0pt}
}

% ============================================================
% COMANDOS MATEMÁTICOS
% ============================================================
\newcommand{\abs}[1]{\left|#1\right|}
\newcommand{\norm}[1]{\left\lVert#1\right\rVert}
\newcommand{\vect}[1]{\mathbf{#1}}
\newcommand{\deriv}[2]{\frac{d#1}{d#2}}
\newcommand{\pderiv}[2]{\frac{\partial#1}{\partial#2}}

% ============================================================
% ENTORNOS / CAJAS ACADÉMICAS
% ============================================================
\newtcolorbox{definition}[1]{
    enhanced,
    breakable,
    title={Definición: #1},
    fonttitle=\bfseries,
    colback=blue!3,
    colframe=blue!50!black,
    boxrule=0.8pt,
    arc=2mm
}

\newtcolorbox{important}{
    enhanced,
    breakable,
    title={Importante},
    fonttitle=\bfseries,
    colback=yellow!8,
    colframe=orange!70!black,
    boxrule=0.8pt,
    arc=2mm
}

\newtcolorbox{example}[1]{
    enhanced,
    breakable,
    title={Ejemplo: #1},
    fonttitle=\bfseries,
    colback=green!4,
    colframe=green!40!black,
    boxrule=0.8pt,
    arc=2mm
}

\newtcolorbox{formula}[1]{
    enhanced,
    breakable,
    title={Fórmula / Regla: #1},
    fonttitle=\bfseries,
    colback=purple!4,
    colframe=purple!50!black,
    boxrule=0.8pt,
    arc=2mm
}

\newtcolorbox{exercise}[1]{
    enhanced,
    breakable,
    title={Ejercicio: #1},
    fonttitle=\bfseries,
    colback=gray!5,
    colframe=gray!60!black,
    boxrule=0.8pt,
    arc=2mm
}

\newtcolorbox{note}{
    enhanced,
    breakable,
    title={Nota},
    fonttitle=\bfseries,
    colback=white,
    colframe=gray!50,
    boxrule=0.6pt,
    arc=2mm
}

% ============================================================
% DOCUMENTO
% ============================================================
\begin{document}

% ---------------- PORTADA ----------------
\begin{titlepage}

\thispagestyle{empty}

\begin{center}

\vspace*{1cm}

{\large \textsc{\universidad}}

\vspace{3cm}

{\Huge \textbf{\materia}}

\vspace{1cm}

{\LARGE \titulo}

\vspace{0.8cm}

\textit{Comprender el razonamiento detrás de cada escrito es lo que transforma el estudio de la materia en criterio profesional.}

\vspace{1.2cm}

\rule{0.70\textwidth}{0.5pt}

\vspace{0.5cm}

{\large Apuntes de estudio}

\vspace{0.5cm}

\rule{0.70\textwidth}{0.5pt}

\vfill

{\large \autor}

\vspace{0.3cm}

{\large \anio}

\end{center}

\vfill

\begin{flushleft}
\textit{Este es un modesto aporte a los estudiantes de ``Derecho Procesal Civil'' no pretende sustituir el material oficial de la materia.}
\end{flushleft}

\end{titlepage}

% ---------------- ÍNDICE ----------------
\tableofcontents
\clearpage

% ============================================================
% CAPÍTULO 1
% ============================================================
\chapter{Metodología previa a la redacción de la contestación de demanda}

Este bloque corresponde a la clase práctica virtual dedicada a la contestación de demanda y a las excepciones previas, desarrollada dentro del ciclo de clases prácticas del mes de mayo, correspondiente a la cátedra del doctor Jorge Rojas. A lo largo de la exposición se utiliza como caso testigo un accidente de tránsito ocurrido en la intersección de la Avenida Callao y la calle Tucumán, en el cual la señora Díaz, en carácter de actora, reclama daños al demandado. Este caso recorre transversalmente todos los bloques temáticos de la clase.

\section{Contexto y diferencias metodológicas con la confección de la demanda}

La confección del escrito de contestación de demanda presenta una lógica de trabajo distinta a la de la demanda. Mientras que, al redactar la demanda, el abogado parte del diálogo mantenido con el cliente y de la documentación y prueba disponibles, al contestar la demanda el profesional ya cuenta con el caso planteado por la parte actora y debe trabajar necesariamente sobre esa base.

\begin{important}
El trabajo de contestación de la demanda se encuentra condicionado por la forma en que la parte actora confeccionó su propio relato. A diferencia de la demanda, donde el relato se construye desde cero, en la contestación la estrategia y la narrativa del demandado deben elaborarse en función de lo ya expuesto por el actor.
\end{important}

\section{Metodología de trabajo: construcción del relato del demandado}

Antes de redactar el escrito, se recomienda seguir un proceso metódico de análisis, compuesto por los siguientes pasos:

\begin{enumerate}
    \item \textbf{Esquematización de los hechos del actor.} Se leen los hechos relatados en la demanda y se anotan uno por uno en una hoja de trabajo.
    \item \textbf{Selección y descarte dentro del relato propio.} Al construir la versión del demandado se descartan dos tipos de hechos: (a) los que no resultan conducentes al caso, y (b) los que, aun siendo conducentes, incriminan al cliente.
    \item \textbf{Confrontación de los dos relatos.} Se colocan en paralelo el relato del actor y el del demandado, a fin de identificar puntos de coincidencia y de divergencia.
    \item \textbf{Análisis de la prueba del actor.} Se estudia qué prueba ofreció la parte actora y se evalúa cuáles de los hechos que a ella le corresponde probar podrían efectivamente acreditarse en la práctica.
    \item \textbf{Análisis de la prueba propia.} Se analizan las pruebas con las que cuenta el demandado.
    \item \textbf{Incorporación selectiva al relato.} Con cautela, se incorporan al relato del actor aquellos elementos del relato propio que resulten útiles para la defensa.
    \item \textbf{Encuadre jurídico.} Más allá del encuadre legal dado por el actor a su demanda, el demandado debe encuadrar también el conflicto dentro del marco normativo correspondiente, a fin de determinar qué disposiciones lo rigen y cómo lo trata la jurisprudencia.
\end{enumerate}

\begin{note}
El relato del demandado no necesariamente coincidirá con el del actor, pero su construcción siempre debe partir de aquel: constituye, en términos prácticos, una nueva versión de lo relatado por la contraparte. Una vez reunido este material —descripto como la ``materia prima''— recién entonces resulta posible redactar el escrito de contestación de demanda.
\end{note}

\section{Consideraciones generales de redacción de los hechos}

Se enuncian una serie de reglas de estilo aplicables a la narración de los hechos dentro del escrito:

\begin{itemize}
    \item No ir y venir en el tiempo, ya que ello genera confusión y dificulta la comprensión del relato; es preferible seguir un esquema cronológico.
    \item La línea argumental debe mantenerse ininterrumpida.
    \item Debe evitarse toda disquisición que aparte la atención del foco del relato.
    \item Los hechos deben narrarse de la forma más sintética y clara posible.
\end{itemize}

\begin{important}
La contestación de demanda, al igual que la demanda misma, no se construye sentándose y escribiendo sin más. No se trata de una narración de ficción, sino de un relato referido a hechos y personas reales, sobre los cuales debe tomarse previamente conocimiento acabado. No corresponde tomar un modelo preelaborado y aplicarlo mecánicamente al caso: es necesario examinar el conflicto concreto y construir una estrategia, a la que luego sigue una táctica para desarrollarla.
\end{important}

\section{Estructura general del escrito}

Para confeccionar el escrito debe conocerse, en primer lugar, la estructura genérica de la contestación de demanda —sus partes y el orden en que se suceden— y, a partir de ella, integrar el desarrollo de cada parte organizado en capítulos, distribuidos según un orden lógico y coherente.

\section{Tasas de justicia y bono del Colegio de Abogados}

A diferencia de lo que ocurre al iniciar la demanda —oportunidad en la que corresponde pagar la tasa de justicia o, en su defecto, iniciar el beneficio de litigar sin gastos—, la contestación de demanda no está sujeta al pago de la tasa de justicia. El demandado puede iniciar un beneficio de litigar sin gastos, pero no está obligado a hacerlo conjuntamente con la contestación.

\begin{important}
\textbf{Obligación fiscal al contestar la demanda — Ley 23.187.} Al momento de presentar la contestación de demanda debe cumplirse con lo dispuesto por la Ley 23.187, al igual que al momento de interponer la demanda, acompañando el bono del Colegio de Abogados. Sin embargo, la contestación de demanda \textbf{no} está sujeta al pago de la tasa de justicia ni a la iniciación del beneficio de litigar sin gastos como requisito para su presentación.
\end{important}

% ============================================================
% CAPÍTULO 2
% ============================================================
\chapter{Estructura del escrito: capítulo de hechos y capítulo de rubros indemnizatorios}

\section{Capítulo de los hechos (artículo 356 del CPCCN)}

Cada articulación defensiva requiere un fundamento de hecho, que se desarrolla en el capítulo del escrito referido a los hechos.

\begin{formula}{Artículo 356 del CPCCN}
El artículo 356 del Código Procesal Civil y Comercial de la Nación le exige al demandado lo mismo que el artículo 330 le exige al actor: una explicación clara y específica de los hechos. El demandado, al igual que el actor en la demanda, realiza una preselección de los hechos que le convienen.
\end{formula}

Frente al relato del actor, el demandado cuenta con las siguientes opciones:

\begin{enumerate}
    \item Negar los hechos.
    \item Contraponer defensas y tensiones.
    \item Contradecir la forma y el alcance de los sucesos.
\end{enumerate}

Además, el demandado puede incorporar hechos que el actor no incluyó en su relato. Lo más probable es que el actor no los haya incorporado porque no le convenían, aunque también puede tratarse de una simple omisión. Para ello, resulta necesario recurrir al borrador elaborado oportunamente e incorporar los hechos que el cliente aportó, de modo que el relato final, al momento del dictado de la sentencia, contemple también los hechos aportados por el demandado.

\subsection{Ejemplo práctico: hechos del actor y contraposición del demandado}

Sobre la base del caso de autos (accidente de tránsito en Av. Callao y Tucumán), puede ilustrarse la lógica de confrontación entre el relato del actor y la posición del demandado:

\begin{table}[H]
\centering
\caption{Confrontación entre el relato del actor y la posición del demandado}
\begin{tabularx}{\textwidth}{>{\raggedright\arraybackslash}X >{\raggedright\arraybackslash}X}
\toprule
\textbf{Hechos del actor (demanda)} & \textbf{Posición del demandado} \\
\midrule
El accidente se produjo el día y hora indicados en la demanda. & El demandado no niega la fecha ni la hora del accidente. \\
El actor estaba detenido en el semáforo. & El demandado se opone cuando el actor afirma que él circulaba a alta velocidad. \\
El demandado circulaba a alta velocidad al momento de la colisión. & El demandado niega haber circulado a alta velocidad. \\
\bottomrule
\end{tabularx}
\end{table}

Además de negar o contraponer los hechos alegados por el actor, el demandado puede incorporar hechos propios, no contradictorios con los del actor pero no incluidos en la demanda: en el caso de autos, que el demandado se encontraba detenido detrás del actor en el mismo semáforo, y que fue el actor quien, al dar marcha atrás mediante una maniobra imprudente, colisionó contra la parte delantera de su vehículo.

\begin{note}
Estos hechos se denominan ``hechos no novedosos'', categoría distinta de los hechos nuevos previstos en el artículo 365 del CPCCN. Se trata de hechos propios del demandado que simplemente no figuran en la demanda. Conviene incorporarlos en un capítulo especial de la contestación, que puede denominarse ``La realidad de los hechos'' o ``La verdad de los hechos''. Mediante este capítulo no siempre podrán desvirtuarse en su totalidad los dichos del actor, pero al menos podrá atenuarse su efecto, lo que en la práctica suele traducirse en una reducción de los rubros reclamados.
\end{note}

\section{Capítulo de impugnación de los rubros indemnizatorios}

El actor cuantifica sus reclamos en distintos rubros indemnizatorios. A modo de ejemplo se mencionan: el daño emergente, el lucro cesante, los gastos de forma (traslados médicos, entre otros) y la desvalorización del vehículo. Corresponde impugnar dichos rubros en forma individual, a los fines de poder impugnar posteriormente el monto final reclamado en la demanda.

% ============================================================
% CAPÍTULO 3
% ============================================================
\chapter{Redacción, forma y trámite del escrito}

\section{Actualización monetaria}

Cuando la demanda contiene un pedido de actualización monetaria, la contestación debe incluir la crítica a la procedencia de esa actualización y/o a la forma en que se practica. Para ello, debe exponerse el modo correcto de actualización y el resultado al que se arriba a través de dicho modo, a fin de demostrar que, más allá del monto reclamado por el actor, este último aplica un método equivocado para actualizarlo.

\section{Lenguaje categórico}

El lenguaje categórico constituye una regla común a la demanda, a la contestación y a cualquier otro planteo que se formule durante el desarrollo del proceso.

\begin{table}[H]
\centering
\caption{Expresiones a evitar y sus equivalentes categóricos}
\begin{tabularx}{\textwidth}{>{\raggedright\arraybackslash}X >{\raggedright\arraybackslash}X}
\toprule
\textbf{Expresión incorrecta (evitar)} & \textbf{Expresión correcta (usar)} \\
\midrule
``Se puede observar\ldots'' & ``Se observa\ldots'' \\
``Se puede concluir\ldots'' & ``Se concluye\ldots'' \\
``No estaría probado\ldots'' & ``No se probó\ldots'' \\
\bottomrule
\end{tabularx}
\end{table}

\begin{important}
Si no se emplea un lenguaje categórico, no podrá lograrse aquello que constituye el objetivo de quien litiga: generar convicción o, cuando menos, una duda razonable. De allí la necesidad de ser categórico y firme en las afirmaciones. Asimismo, no se argumenta mejor por repetir varias veces lo mismo ni por emplear más palabras de las necesarias para expresar una idea que podría formularse de manera más sintética.
\end{important}

\section{Citas normativas y jurisprudenciales}

Corresponde transcribir literalmente las normas legales cuando:

\begin{itemize}
    \item Son muy específicas desde el punto de vista técnico.
    \item Se trata de leyes secretas, aún no desclasificadas.
    \item La cuestión regulada es oscura o compleja.
\end{itemize}

En cambio, no corresponde transcribir literalmente las normas de uso común y corriente. La cita de las normas debe ser lo más idéntica y sintética posible, evitando todo desarrollo innecesario.

En cuanto a las citas jurisprudenciales, resultan muy ilustrativas para el tribunal, en especial los fallos de las Cortes —privilegiando, dentro de estos, los más recientes— y, en segundo orden, los fallos de la Cámara de Apelaciones del fuero en el que se litiga.

\section{Trámite del escrito en soporte papel}

El escrito de contestación de demanda debe entregarse en soporte papel siguiendo el siguiente procedimiento:

\begin{enumerate}
    \item Se presenta el original en la secretaría del juzgado, dentro del plazo para contestar la demanda.
    \item Junto con el escrito original se acompaña un juego completo de copias de toda la prueba instrumental, que se incorpora al expediente.
    \item Se presenta una copia simple del escrito —solo del escrito— para que sea sellada como constancia de presentación, la cual el abogado conserva en su carpeta.
    \item Se arma una carpeta paralela al expediente dentro del estudio jurídico, con un juego completo de la documentación.
\end{enumerate}

Respecto de las firmas:

\begin{itemize}
    \item El escrito que se entrega es firmado por la parte y el letrado patrocinante, o solo por el apoderado.
    \item Las copias que se entregan son firmadas únicamente por los letrados; no es necesario que las firme la parte.
    \item Las copias de la documental también son firmadas por el letrado.
\end{itemize}

\begin{note}
Al escrito que se deja en secretaría debe reponérsele el cargo en presencia del letrado. Asimismo, el escrito puede dejarse en la mesa receptora habilitada en Lavalle 1220, planta baja, que funciona para los juzgados civiles radicados en Av. de los Inmigrantes 1950, y para los de Lavalle 1212 y Lavalle 1250. La constancia entregada como recibo lleva el cargo o sello del tribunal o de la mesa receptora y la fecha, lo que acredita la presentación dentro del plazo legal. Si el escrito se presenta dentro de las dos primeras horas del horario judicial, debe hacerse constar además la hora de presentación, la cual debe coincidir con la del cargo mecánico.
\end{note}

\section{Carga digital en el expediente electrónico}

\begin{formula}{Acordada N.\textsuperscript{o} 3/2015 — Poder Judicial de la Nación (punto 5)}
Dentro de las veinticuatro (24) horas posteriores a la entrega del escrito en el juzgado, debe subirse una copia digital en formato PDF al expediente electrónico, a través de la página del Poder Judicial, ingresando con el número de CUIT del letrado. El PDF debe contener, pieza por pieza y en el mismo orden cronológico, todo lo acompañado en papel: el escrito de contestación de demanda y la documental. El incumplimiento de esta obligación puede dar lugar a la intimación del letrado y, como sanción grave, a que se tenga por no presentado el escrito.
\end{formula}

De no cumplirse con la carga del PDF, corresponde una intimación al letrado bajo apercibimiento, siendo la sanción por incumplimiento particularmente grave: tener por no presentado el escrito. Se recomienda imprimir la constancia de visualización del PDF ingresado y conservarla junto con la contestación sellada.

% ============================================================
% CAPÍTULO 4
% ============================================================
\chapter{Articulación de excepciones previas}

Este bloque temático fue desarrollado por un segundo docente, colaborador de la cátedra en la comisión a cargo del doctor Salgado, como continuidad pedagógica de la clase práctica de contestación de demanda ya expuesta.
% TODO: contenido incompleto o ambiguo en las notas originales — la referencia a "doctor Notas" en la transcripción original no resulta clara; se conserva la referencia al docente colaborador de la comisión del doctor Salgado, dato que sí figura de forma expresa en el apunte.

\section{Clasificación: impedimentos procesales, excepciones propiamente dichas y defensas previas}

Corresponde distinguir las defensas nominadas —aquellas que cuentan con un \textit{nomen iuris} otorgado por el ordenamiento de fondo, como la excepción de incumplimiento— de las excepciones previas propiamente dichas. Dentro de estas últimas se distinguen tres categorías:

\begin{table}[H]
\centering
\caption{Categorías de excepciones previas}
\begin{tabularx}{\textwidth}{>{\raggedright\arraybackslash}p{0.22\textwidth} >{\raggedright\arraybackslash}p{0.36\textwidth} >{\raggedright\arraybackslash}X}
\toprule
\textbf{Categoría} & \textbf{A qué apunta} & \textbf{Ejemplo} \\
\midrule
Impedimentos procesales & Cuestionan la constitución válida del proceso (relación procesal). & Incompetencia del juez. \\
Excepciones propiamente dichas & Atacan el ejercicio de la acción (pretensión procesal), no el derecho de fondo. & Prescripción: el derecho existe, pero la acción ya no puede promoverse por haber vencido el plazo. \\
Defensas previas & Atacan directamente el derecho de fondo que funda el reclamo. & Desistimiento del derecho: la parte ya renunció a ese derecho, por lo que no puede volver a promover juicio. \\
\bottomrule
\end{tabularx}
\end{table}

\section{Contexto: el proceso ya se encuentra iniciado}

Al momento de contestar la demanda, el proceso ya ha sido promovido contra el cliente. Desde la carátula del expediente puede obtenerse información relevante para evaluar la procedencia de eventuales excepciones:

\begin{itemize}
    \item El nombre del juez, que brinda información sobre la competencia.
    \item Los nombres de las partes, que pueden sugerir cuestiones de legitimación.
    \item El número de expediente y el juzgado actuante.
\end{itemize}

\section{Análisis inicial de la cédula de notificación}

Cuando el cliente recibe la cédula y la presenta en el estudio jurídico, el letrado debe examinarla detenidamente, dado que pueden existir cuestiones que ameriten planteos previos:

\begin{enumerate}[label=\Alph*.]
    \item Defectos en la cédula misma, que podrían habilitar un planteo de nulidad procesal.
    \item Fecha en que fue ordenado el traslado de la demanda, a fin de verificar si el expediente tuvo impulso útil dentro de los plazos legales, lo que permitiría evaluar una eventual caducidad de instancia.
\end{enumerate}

\begin{important}
\textbf{Plazo diferenciado para nulidad y caducidad (CPCCN).} Tanto el planteo de nulidad como el de caducidad de instancia cuentan con un plazo de cinco (5) días para ser presentados, contados desde la notificación. Si el planteo no se deduce dentro de ese plazo, el vicio queda purgado y la falta de impulso queda curada. Este plazo es diferente e independiente del plazo de quince (15) días hábiles para contestar la demanda en el proceso ordinario (o el que resulte extendido por la distancia, según el artículo 158 del CPCCN). Constituye un error frecuente en la práctica intentar plantear la nulidad el día 15, junto con la contestación, cuando el plazo para hacerlo ya venció el día 5.
\end{important}

\section{Oportunidad para articular las excepciones previas}

Las excepciones previas se deducen junto con el escrito de contestación de demanda, es decir, en la misma presentación. No obstante, existen tres excepciones que suspenden el plazo para contestar la demanda:

\begin{enumerate}
    \item Excepción de defecto legal.
    \item Excepción de arraigo.
    \item Excepción de falta de personalidad.
\end{enumerate}

La excepción de defecto legal, en su carácter de suspensiva del plazo para contestar la demanda, responde a una política razonable, en tanto hace al ejercicio del derecho de defensa en juicio: si la demanda resulta incomprensible o la parte demandada no puede comprender los alcances de la pretensión, se pone en crisis dicho derecho. Respecto de la excepción de arraigo, en la práctica se presenta con escasa frecuencia, dado que con el Código Civil y Comercial de la Nación su alcance quedó fuertemente limitado.

\section{Ubicación de las excepciones dentro del escrito y orden de prelación}

El planteo de las excepciones dentro del cuerpo del escrito de contestación de demanda se sitúa en forma anterior a la contestación propiamente dicha: es decir, antes de la negativa, del eventual reconocimiento y del eventual relato que la parte quiera efectuar. Este orden resulta lógico, en la medida en que el juez debe resolver las excepciones en forma previa al tratamiento del asunto de fondo.

La hoja de ruta general del escrito es, en consecuencia, la siguiente:

\begin{enumerate}
    \item Sumario del escrito.
    \item Presentación de quien comparece.
    \item Acreditación de la personería.
    \item Planteos de nulidad o de caducidad de instancia, si correspondieren (plazo: 5 días).
    \item Excepciones previas, si se articulan.
    \item Contestación de demanda propiamente dicha (negativa, reconocimiento, relato, rubros).
\end{enumerate}

No existe una norma específica que regule el orden entre las distintas excepciones previas; sin embargo, en los usos forenses se suele articular en primer término los impedimentos procesales, para luego pasar a las excepciones que hacen a la relación jurídica de fondo.

\begin{formula}{CPCCN — Orden de resolución de las excepciones}
El Código Procesal Civil y Comercial de la Nación establece que el juez debe resolver primero las excepciones de competencia y de litispendencia, antes de analizar las demás. Este criterio responde a razones lógicas: si el juez declara su propia incompetencia, ya no puede inmiscuirse en el tratamiento de los demás asuntos.
\end{formula}

\section{Ejemplo de escrito real: estructura y prueba de las excepciones}

A partir de un ejemplo de escrito real en el que se contesta la demanda y se oponen excepciones, pueden identificarse los siguientes lineamientos, coincidentes con lo ya expuesto:

\begin{itemize}
    \item Título del escrito y sumario.
    \item Identificación de quien se presenta (nombre, domicilio, letrado que lo representa).
    \item Acreditación de la personería.
    \item Articulación de las excepciones —en el ejemplo examinado: falta de personería y litispendencia—.
\end{itemize}

\begin{important}
Debe distinguirse entre el ofrecimiento de prueba relativo a la contestación de demanda propiamente dicha —esto es, las cuestiones que se pretenden probar respecto del fondo del asunto— y la prueba relativa a la excepción planteada. Existen excepciones que, de no ofrecerse prueba a su respecto, carecen de todo asidero. Así, por ejemplo, la excepción de litispendencia debe estar acompañada necesariamente de un testimonio de la otra causa o, cuando menos, del ofrecimiento de esa otra causa como prueba.
\end{important}

\section{Estrategia procesal: qué excepciones plantear y cuáles no}

No toda excepción que formalmente asista al demandado resulta conveniente de plantear dentro de la estrategia adoptada en el caso concreto.

\begin{example}{Excepción de incompetencia}
Puede presentarse un caso en el que el juez interviniente resulte incompetente, correspondiendo que intervenga, en cambio, un juez con asiento en otra jurisdicción —por ejemplo, en la Provincia de Buenos Aires en lugar de la Ciudad—. Si bien este planteo sería formalmente correcto, el abogado debe formularse ciertas preguntas estratégicas antes de articularlo: si no se encuentra matriculado en la jurisdicción que resultaría competente, no podrá ejercer allí la defensa; asimismo, puede ocurrir que en el fuero que resultaría competente las indemnizaciones sean, en la práctica, normalmente más elevadas que en el fuero de origen, lo que desaconsejaría remitir la causa.
\end{example}

\begin{example}{Excepción de defecto legal}
La excepción de defecto legal presenta utilidad desde dos perspectivas. Por un lado, puede servir como herramienta para lograr la suspensión del plazo para contestar la demanda. Por otro lado, está destinada a remediar una situación en la que el demandado no puede ejercer debidamente su derecho de defensa en juicio: si existen falencias en la demanda —por ejemplo, que el actor no individualizó correctamente los daños sufridos, o no indicó dónde trabajaba ni qué lesión padeció—, tales falencias podrían ser interpretadas en contra de los intereses del actor al momento de dictarse sentencia. En consecuencia, plantear la excepción de defecto legal puede darle al actor la oportunidad de sanear una demanda mal confeccionada.
\end{example}

\begin{important}
No corresponde articular todas las excepciones disponibles por el solo hecho de contar con ellas: la acumulación de excepciones puede resultar contraproducente. Plantear un número excesivo de excepciones no necesariamente beneficia a la parte demandada, en la medida en que algunos planteos pueden resultar contradictorios entre sí y debilitar la postura procesal adoptada. Un ejemplo característico de esta situación, frecuente en los procesos ejecutivos, es la articulación simultánea de la excepción de falsedad de título —sosteniendo que la firma inserta en el documento no pertenece al demandado— junto con la excepción de pago —sosteniendo que la obligación ya fue abonada—. Dicha contradicción debilita fuertemente la postura de quien se defiende e implica la desestimación, sin mayor análisis, de una de ambas defensas.
\end{important}

% ============================================================
% CUADRO CORNELL
% ============================================================
\chapter{Cuadro Cornell de repaso}

El siguiente cuadro reúne, según el método Cornell, las preguntas disparadoras y las respuestas clave correspondientes a los cuatro bloques temáticos desarrollados: la metodología previa a la redacción, la estructura del escrito, las consideraciones de forma y trámite, y la articulación de excepciones previas. Se incluye, al final, una síntesis integradora.

\begin{longtable}{
    >{\raggedright\arraybackslash}p{0.32\textwidth}
    >{\raggedright\arraybackslash}p{0.58\textwidth}
}
\caption{Cuadro Cornell — Contestación de demanda y excepciones previas} \\
\toprule
\textbf{Preguntas / palabras clave} & \textbf{Notas / respuestas} \\
\midrule
\endfirsthead

\toprule
\textbf{Preguntas / palabras clave} & \textbf{Notas / respuestas} \\
\midrule
\endhead

\midrule
\multicolumn{2}{r}{\textit{Continúa en la página siguiente}} \\
\endfoot

\bottomrule
\endlastfoot

\textbf{¿En qué se diferencia metodológicamente contestar una demanda de confeccionarla?} &
Al redactar la demanda, el abogado parte de un diálogo propio con el cliente. Al contestarla, el trabajo está condicionado por el relato construido por el actor: debe partirse de esa demanda para elaborar una versión propia, seleccionando los hechos convenientes y descartando los que incriminan al cliente. \\

\textbf{¿Cuáles son los pasos previos a la redacción del escrito?} &
1. Anotar todos los hechos de la demanda. 2. Seleccionar y descartar. 3. Confrontar los dos relatos. 4. Analizar la prueba del actor. 5. Analizar la prueba propia. 6. Incorporar selectivamente. 7. Encuadrar jurídicamente. \\

\textbf{¿Qué puede hacer el demandado frente a los hechos del actor?} &
Negarlos, contraponer defensas y tensiones, contradecir su forma y alcance, e incorporar hechos omitidos por el actor que resulten convenientes para el demandado. \\

\textbf{¿Qué exige el artículo 356 del CPCCN al demandado?} &
Una explicación clara y específica de los hechos, al igual que el artículo 330 se lo exige al actor. El demandado también realiza una preselección de los hechos que le convienen. \\

\textbf{¿Qué son los ``hechos no novedosos'' del demandado?} &
Son hechos propios del demandado que no figuran en la demanda, categoría distinta de los ``hechos nuevos'' del art. 365 del CPCCN. Ejemplo: el demandado estaba detenido detrás del actor y fue este quien, al dar marcha atrás, chocó su auto. Se incorporan en un capítulo denominado ``La verdad de los hechos''. \\

\textbf{¿Debe el demandado pagar tasa de justicia al contestar?} &
No. La contestación de demanda no está sujeta al pago de la tasa de justicia. Sí debe cumplirse con la Ley 23.187 y acompañar el bono del Colegio de Abogados. El beneficio de litigar sin gastos puede iniciarse, pero no es obligatorio hacerlo al momento de contestar. \\

\textbf{¿Qué plazo existe para plantear nulidad o caducidad de instancia?} &
Cinco días desde la notificación. Este plazo es diferente e independiente del plazo de quince días para contestar la demanda. Si no se deduce en término, el vicio queda purgado; constituye un error frecuente pretender plantearlo el día 15. \\

\textbf{¿Cuándo se articulan las excepciones previas?} &
Junto con el escrito de contestación de demanda, en la misma presentación. Sin embargo, tres excepciones suspenden el plazo para contestar: defecto legal, falta de personalidad y arraigo. \\

\textbf{¿Dónde se ubican las excepciones dentro del escrito?} &
Antes de la contestación de demanda propiamente dicha: antes de la negativa, el reconocimiento y el relato. Primero se articulan los impedimentos procesales y luego las excepciones que hacen a la relación de fondo. \\

\textbf{¿Qué prueba debe ofrecerse al articular excepciones?} &
Cada excepción debe estar sustentada en prueba. Por ejemplo, la de litispendencia debe acompañarse necesariamente de un testimonio de la otra causa, o al menos del ofrecimiento de esa otra causa como prueba; sin prueba, la excepción carece de asidero. \\

\textbf{¿Por qué no conviene acumular demasiadas excepciones?} &
Porque pueden resultar contradictorias entre sí y debilitar la postura del demandado. Ejemplo clásico: plantear falsedad de firma y excepción de pago en la misma presentación resulta contradictorio e implica la desestimación de una de ellas. \\

\textbf{¿Cuándo conviene no articular una excepción de incompetencia?} &
Cuando el fuero competente sea aquel en el que el abogado no se encuentra matriculado, o cuando en el fuero donde se inició el proceso las indemnizaciones suelan ser más elevadas que en el que resultaría competente. \\

\textbf{¿Cuándo debe subirse el PDF al expediente electrónico?} &
Dentro de las 24 horas posteriores a la entrega del escrito en papel. El PDF debe reflejar exactamente, y en el mismo orden cronológico, todo lo presentado en papel. La obligación surge del punto 5 de la Acordada 3/2015; su incumplimiento puede dar lugar a que se tenga el escrito por no presentado. \\

\textbf{¿Qué reglas de estilo rigen la redacción del escrito?} &
Ser categórico (``se observa'', no ``se puede observar''); seguir un orden cronológico; evitar disquisiciones; ser sintético y claro; citar jurisprudencia de la Corte (fallos recientes) y de la Cámara del fuero; transcribir las normas literalmente solo cuando sean muy técnicas u oscuras. \\

\midrule
\multicolumn{2}{p{0.92\textwidth}}{
\textbf{Resumen:} La contestación de demanda es un acto procesal estratégico que exige un trabajo previo riguroso: analizar la demanda del actor, construir una versión propia de los hechos seleccionando y descartando con criterio, confrontar ambos relatos y encuadrar jurídicamente el conflicto. El escrito posee una estructura capitular definida (sumario, personería, excepciones, hechos, rubros, prueba) y debe redactarse con lenguaje categórico, cronológico y sintético. Desde el punto de vista fiscal, no requiere tasa de justicia, pero sí el bono del Colegio de Abogados y el cumplimiento de la Ley 23.187. El trámite exige presentación en papel en secretaría y carga digital en el expediente electrónico dentro de las 24 horas. Respecto de las excepciones previas, deben articularse en el mismo escrito —antes de la contestación propiamente dicha—, sustentadas en prueba y en un orden que priorice los impedimentos procesales. La estrategia procesal resulta determinante: existen excepciones formalmente procedentes que pueden ser perjudiciales para el cliente, como revelar deficiencias de la demanda al actor o remitir la causa a un fuero desfavorable. La acumulación irreflexiva de excepciones contradictorias debilita la postura del demandado. En síntesis, contestar una demanda no consiste en completar un formulario, sino en construir, con método y estrategia, la defensa integral del cliente.
} \\

\end{longtable}

\end{document}
